\chapter{Crush, drama và đầu óc hay mất sóng}
\markboth{Crush, drama và đầu óc hay mất sóng}{Crush, drama và đầu óc hay mất sóng}

\section{Em thất tình nên đầu em treo máy}
\begin{truyen}{Em thất tình nên đầu em treo máy}{Classroom Talk}
\chuhoa{M}{ột bạn ngồi thừ ra rồi thở dài: ``Thầy thông cảm, em đang thất tình nên đầu em treo máy mấy bữa nay.''}

Thầy gật: ``Đau lòng là thật. Nhưng em định để một người rời khỏi đời mình tiện tay dắt luôn lịch học, nề nếp và tương lai đi theo à?''

Bạn ấy cắn môi: ``Nghe hơi phũ nha thầy.''

Thầy đáp: ``Vì nỗi buồn cần được ôm, nhưng không nên được trao luôn quyền điều hành cả cuộc đời em.''

\textit{(Heartbreak is real. Letting it run your entire life is a second mistake on top of the first wound.)}
\end{truyen}

\begin{baihoc}
Đau vì một người là chuyện người. Tự bỏ mình theo nỗi đau đó mới là chuyện dại.
\textit{(Being hurt by someone is human. Abandoning yourself because of that hurt is the foolish part.)}
\end{baihoc}

\begin{ghinhoanh}
\textbf{Short English to remember:}

Feel the pain.

Do not hand it the steering wheel.
\end{ghinhoanh}

\begin{tuhocnhanh}
1. Em đang buồn thật, hay đang nuôi nỗi buồn để hợp thức hóa việc buông xuôi? \textit{(Are you truly grieving, or using sadness to justify giving up?)}

2. Một việc nhỏ nào em vẫn có thể giữ cho mình trong giai đoạn này? \textit{(What small routine can you still protect during this period?)}
\end{tuhocnhanh}
\ngancach

\section{Ngồi gần crush sao tập trung nổi}
\begin{truyen}{Ngồi gần crush sao tập trung nổi}{Classroom Talk}
\chuhoa{C}{ó bạn than rất chân thành: ``Thầy cho em đổi chỗ được không? Em ngồi gần bạn đó là đầu em mất hết tín hiệu học tập.''}

Thầy hỏi: ``Một gương mặt mà làm bay sạch phương trình, định nghĩa và cả trí nhớ ngắn hạn của em, vậy nền học của em mỏng cỡ giấy ăn hả?''

Bạn ấy đỏ mặt: ``Tại em chưa quen thôi.''

Thầy cười: ``Vậy tập quen đi. Chứ sau này ra đời đâu phải ai đẹp là em được quyền tắt não.''

\textit{(Attraction is normal. But if it erases all focus, the issue is bigger than the person in front of you.)}
\end{truyen}

\begin{baihoc}
Cảm xúc mạnh không đáng xấu hổ. Khả năng tự lái mình yếu mới đáng lo.
\textit{(Strong feelings are not shameful. Weak self-control is the deeper concern.)}
\end{baihoc}

\begin{ghinhoanh}
\textbf{Short English to remember:}

Attraction is normal.

Self-control is still required.
\end{ghinhoanh}

\begin{tuhocnhanh}
1. Em đang rung động, hay đang lấy rung động làm cớ cho việc mất tập trung? \textit{(Are you attracted, or using attraction as an excuse for losing focus?)}

2. Em cần luyện kỹ năng gì để cảm xúc không đá mình khỏi đường ray? \textit{(What skill do you need so emotions stop kicking you off track?)}
\end{tuhocnhanh}
\ngancach

\section{Bạn ấy seen không rep nên em tụt mood cả tuần}
\begin{truyen}{Bạn ấy seen không rep nên em tụt mood cả tuần}{Classroom Talk}
\chuhoa{M}{ột bạn buồn rười rượi: ``Người ta seen không rep có một cái mà em tụt mood cả tuần, học gì nổi nữa thầy.''}

Thầy hỏi: ``Một thông báo im lặng mà điều khiển được cả bảy ngày của em, vậy cuộc đời này em đang cầm hay đang cho người khác giữ hộ?''

Bạn ấy cười méo: ``Nghe nhục ghê.''

Thầy gật: ``Nhục một chút còn hơn sống kiểu ai cũng có quyền chạm màn hình rồi điều chỉnh tâm trạng em.''

\textit{(A missed reply can sting, but giving it control over an entire week is an expensive surrender.)}
\end{truyen}

\begin{baihoc}
Đừng giao nguyên tuần sống của mình cho một người còn không chịu gõ nổi một dòng tử tế.
\textit{(Do not hand over a whole week of your life to someone who cannot even type one decent reply.)}
\end{baihoc}

\begin{ghinhoanh}
\textbf{Short English to remember:}

A silent reply can hurt.

It should not own your week.
\end{ghinhoanh}

\begin{tuhocnhanh}
1. Em đang buồn vì người đó, hay vì giá trị bản thân em đang treo trên phản hồi của họ? \textit{(Are you sad because of that person, or because your self-worth hangs on their response?)}

2. Em cần dựng lại điều gì trong mình để không bị dắt bởi thông báo? \textit{(What do you need to rebuild so notifications stop controlling you?)}
\end{tuhocnhanh}
\ngancach

\section{Yêu vào em mới có động lực học}
\begin{truyen}{Yêu vào em mới có động lực học}{Classroom Talk}
\chuhoa{C}{ó bạn nói rất tự tin: ``Thật ra yêu vào em mới có động lực học đó thầy. Có người để phấn đấu nghe nó máu hơn.''}

Thầy hỏi: ``Nếu động lực em phải cắm điện từ một người khác, lỡ mai họ rút phích thì đời em tối thui luôn à?''

Bạn ấy cười: ``Thầy ví đau quá.''

Thầy đáp: ``Vì mượn người khác làm ổ cắm ý chí là kiểu rất lãng mạn nhưng cực kỳ kém bền.''

\textit{(Love can inspire, but using another person as your only power source makes your discipline dangerously fragile.)}
\end{truyen}

\begin{baihoc}
Người khác có thể tiếp lửa. Nhưng bộ máy phải là của em, không thì cháy hụt rất nhanh.
\textit{(Someone else may light the fire. The engine still has to be yours, or the flame dies quickly.)}
\end{baihoc}

\begin{ghinhoanh}
\textbf{Short English to remember:}

Inspiration may come from others.

Discipline must belong to you.
\end{ghinhoanh}

\begin{tuhocnhanh}
1. Nếu không có người đó, em còn lý do nào để cố không? \textit{(If that person disappears, do you still have a reason to keep going?)}

2. Em đang yêu để lớn lên, hay đang mượn yêu để đỡ phải tự lái mình? \textit{(Are you loving in a way that helps you grow, or using love to avoid steering your own life?)}
\end{tuhocnhanh}
\ngancach